% ==============================================================================
% 02-interfaces.tex — Раздел 2.1. Требования к внешним интерфейсам
% ==============================================================================

\chapter{Детальные требования}
\label{ch:detailed}

\section{Требования к внешним интерфейсам}
\label{sec:interfaces}

\subsection{Интерфейсы пользователя (UI)}
\label{subsec:ui}

Пользовательский интерфейс на базе Next.js, React, Tailwind, shadcn/ui. Экраны: главная страница, вход в~систему, регистрация, чат с~ассистентом (история сообщений, вложения, индикатор выполнения, результаты по платформам), профиль бизнеса, управление подключением каналов (через OAuth для платформ с~программным интерфейсом или по инструкции для браузерной автоматизации), лента отзывов, история публикаций, история задач, настройки учётной записи.

\subsection{Аппаратные интерфейсы}
\label{subsec:hardware}

Клиент~--- современный браузер. Сервер: минимум 2~ядра, 4~ГБ оперативной памяти, 20~ГБ дискового пространства; дополнительно не менее 1~ГБ на агент с~браузерной автоматизацией.

\subsection{Программные интерфейсы}
\label{subsec:software}

Интерфейс REST (аутентификация, бизнес, интеграции, диалоги), аутентификация по токену JWT. Интеграции: VK (OAuth~2.0), Telegram (интерфейс бота), Яндекс.Бизнес (браузерная автоматизация). Документация~--- описание интерфейса по стандарту OpenAPI.

\subsection{Интерфейсы связи и коммуникации}
\label{subsec:communication}

HTTP/HTTPS, потоковая передача ответов чата по HTTP, шина сообщений NATS для протокола «агент--агент». Целевые показатели: 95-й перцентиль времени отклика интерфейса < 500~мс, ответ языковой модели < 10~с.
